\mxcorrections{5}\label{C05-CORR}
\textbf{Usage professeur.} Les réponses qui suivent concernent les tâches non résolues dans le cours. Faire d’abord expliciter la démarche ; fournir une aide graduée avant la solution. Les durées sont des estimations à étalonner en classe.
\subsection*{Diagnostic et remédiation}
\begin{enumerate}
\item $\dfrac{12}{40}=\dfrac{3}{10}=\pc{30}$. \textit{Action :} Relier effectif partiel et effectif de référence.
\item $\num{0,7}$. \textit{Action :} Représenter l’univers partagé en deux événements.
\item $A\cap B$. \textit{Action :} Colorier l’intersection de deux ensembles.
\item $A\cup B$. \textit{Action :} Inclure les issues communes dans la réunion.
\item $\num{0,1}$. \textit{Action :} Remplacer les décimaux par des fractions.
\item $\dfrac{3}{6}=\dfrac{1}{2}$. \textit{Action :} Identifier l’hypothèse d’équiprobabilité.
\end{enumerate}
\subsection*{Activité de recherche}
Intersection nécessaire 12 ; cases 12,28,18,42. Remplacer par 13,27,17,43 garde les marges mais $\dfrac{13}{100}\ne \num{0,4}\times \num{0,3}$.
\subsection*{Applications immédiates}
\begin{enumerate}
\item $\dfrac{5}{60}=\dfrac{1}{12}$ et $\dfrac{5}{20}=\dfrac{1}{4}$.
\item $\dfrac{1}{4}+\dfrac{1}{6}=\dfrac{5}{12}$.
\item Non : $\num{0,4}\times\num{0,5}=\num{0,20}\ne\num{0,18}$.
\item $(\dfrac{2}{3})^4=\dfrac{16}{81}$ puis $1-\dfrac{16}{81}=\dfrac{65}{81}$.
\item $\dfrac{675}{2000}=\num{0,3375}$ ; c’est une estimation, pas la valeur exacte.
\end{enumerate}
\subsection*{Exercices : démarche et résultats}
\paragraph{Exercice 5.1 -- Un tableau.} Cases : 15 (A et B), 25 (A sans B), 20 (B sans A), 40 (ni A ni B). $P_A(B)=\dfrac{15}{40}=\dfrac{3}{8}$ ; $P_B(A)=\dfrac{15}{35}=\dfrac{3}{7}$.
\paragraph{Exercice 5.2 -- Production.} $P(D)=\num{0,25}\times\num{0,04}+\num{0,75}\times\num{0,02}=\num{0,025}$. $P_D(A)=\dfrac{\num{0,01}}{\num{0,025}}=\num{0,4}$. Le dénominateur change lorsqu’on conditionne par D.
\paragraph{Exercice 5.3 -- Indépendance à vérifier.} $P(A)=P(B)=\dfrac{1}{2}$ ; intersection $\ens{4;6}$ de probabilité $\dfrac{1}{3}\ne \dfrac{1}{4}$ : non indépendants. Contrairement à l’exemple du cours (résultat au moins 5), ici le seuil 4 ne donne pas l’indépendance.
\paragraph{Exercice 5.4 -- Sans remise.} Deux rouges : $(\dfrac{3}{5})(\dfrac{2}{4})=\dfrac{3}{10}$. Une de chaque : $(\dfrac{3}{5})(\dfrac{2}{4})+(\dfrac{2}{5})(\dfrac{3}{4})=\dfrac{3}{5}$. L’ordre distingue deux chemins disjoints.
\paragraph{Exercice 5.5 -- Trois essais.} Un succès : trois chemins, chacun $(\dfrac{1}{3})(\dfrac{2}{3})^2=\dfrac{4}{27}$, total $\dfrac{4}{9}$. Deux succès : $3(\dfrac{1}{3})^2(\dfrac{2}{3})=\dfrac{2}{9}$ ; trois : $\dfrac{1}{27}$ ; total $\dfrac{7}{27}$.
\paragraph{Exercice 5.6 -- Détecteur.} $P(S)=\num{0,02}(\num{0,9})+\num{0,98}(\num{0,05})=\num{0,067}$. $P_S(D)=\dfrac{\num{0,018}}{\num{0,067}}=\dfrac{18}{67}\approx \pc{26,9}$. Le faible taux initial explique l’écart avec \pc{90}.
\paragraph{Exercice 5.7 -- Ensembles et cardinal.} $A\cap B=\{2\}$ (1) ; $A\cup B=\ens{1;2;4}$ (3) ; $E\setminus A=\ens{3;4}$ (2) ; produit $\{\pt{1;2},\pt{1;4},\pt{2;2},\pt{2;4}\}$ (4). Un couple est ordonné.
\paragraph{Exercice 5.8 -- Reprise différée.} $P(B)=\dfrac{1}{8}+\dfrac{3}{8}=\dfrac{1}{2}$. $P_B(A)=\dfrac{\dfrac{1}{8}}{\dfrac{1}{2}}=\dfrac{1}{4}$. L’égalité numérique avec $P_A(B)$ est particulière à ces données.
\paragraph{CQFD, exercice 67 p. 336, questions 1, 2 et 3.} $P(\overline D\cap\overline N)=\num{0,9}\times\num{0,05}=\num{0,045}$. La probabilité d’un test positif vaut $\num{0,1}\times\num{0,97}+\num{0,045}=\num{0,142}$. Parmi les tests positifs, la probabilité de ne pas être dopé vaut $\dfrac{\num{0,045}}{\num{0,142}}=\dfrac{45}{142}\approx\pc{31,7}$. Ne pas inverser test positif et négatif.
\paragraph{CQFD, exercice 70 p. 336, questions a et b.} $P(B)=\dfrac14$, $P_A(B)=\dfrac12$ et $P_{\overline A}(B)=0$. Ainsi $\dfrac14\ne\dfrac12+0$, puis $0\ne1-\dfrac12$. La formule des probabilités totales comporte les poids $P(A)$ et $P(\overline A)$ ; le complément s’applique à l’événement $B$ sous une même condition : $P_A(B)+P_A(\overline B)=1$, et non $P_A(B)+P_{\overline A}(B)=1$.
\subsection*{QCM : réponses et raisons}
\textbf{QCM 1.}
\begin{enumerate}
\item B. Univers de référence A.
\item A. Produit des branches.
\item D. Définition.
\item C. Critère sans division.
\item B. Intersection nulle, produit positif.
\end{enumerate}
\textbf{QCM 2.}
\begin{enumerate}
\item C. Complément de trois faces.
\item A. Produit $p^2$.
\item D. Probabilités d’une partition conditionnelle.
\item B. Fréquence observée et probabilité du modèle diffèrent.
\item C. Il reste 2 rouges sur 4 boules.
\end{enumerate}
\subsection*{Vrai ou faux}
\begin{enumerate}
\item Faux : 12 objets communs sur 40 A et 30 B donnent $\dfrac{3}{10}$ et $\dfrac{2}{5}$.
\item Vrai : $P(A\cap\overline B)=P(A)-P(A\cap B)=P(A)(1-P(B))$.
\item Faux : sans remise, la composition de l’urne change.
\end{enumerate}
\paragraph{Différenciation.} Pour un blocage technique, proposer la première étape seulement (domaine, facteur ou univers de référence), puis retirer l’aide. Pour un élève autonome, demander une seconde méthode et le traitement d’un cas limite. À l’oral, autoriser une reformulation initiale dans la langue maîtrisée, puis exiger une phrase mathématique précise en français.
